\subsection{\texorpdfstring{\(\mathcal{F}_{\mathfrak{S}}(\mathbf{X};\mathbf{Y})\) 的完全子集}{F 下标 S (X;Y) 的完全子集}}

命题 5. 设 \(\Phi\) 是一个集，Y 是一致空间，\(\mathfrak{S}\) 是 X 的子集组成的集。则滤子 \(\Phi\)（\(\mathcal{F}_{\mathfrak{S}}(\mathbf{X};\mathbf{Y})\) 上的）收敛于 \(u_0\) 的充要条件是：\(\Phi\) 是关于 \(\mathfrak{S}\) -收敛的一致结构的柯西滤子，且逐点收敛于 \(u_0\)（在 \(\mathbf{B} = \bigcup_{\mathbf{A}\in \mathfrak{S}}\mathbf{A}\) 中）。

由于 B 中逐点收敛的结构比 S-收敛的结构粗，只需证明：对每个 \(A \in S\) 与 Y 的每个闭围域 V，\(W(A, V)\) 关于逐点收敛拓扑在 B 中是闭的（第 II 章，§ 3，第 3 号，命题 7）。而 \(W(A, V)\) 是诸映射 \((u, v) \to (u(x), v(x))\)（当 x 取遍 A 时）下 V 的原象的交；这些映射关于逐点收敛拓扑连续（第 2 号注记 6），由此即得结论。

推论 1. \(\mathcal{F}_{\mathfrak{S}}(\mathbf{X};\mathbf{Y})\) 的子空间 H 是完全的，充要条件是：对 H 上每个柯西滤子 \(\Phi\)，存在 \(u_0\in \mathbf{H}\)，使 \(\Phi\) 逐点收敛于 \(u_0\)（在 \(\mathbf{B} = \bigcup_{\mathbf{A}\in \mathfrak{S}}\mathbf{A}\) 中）。

这由命题 5 立得。

推论 2. 设 \(\mathfrak{S}_1, \mathfrak{S}_2\) 是 X 的子集组成的两个集，它们的并相同，且 \(\mathfrak{S}_1 \subset \mathfrak{S}_2\)，又设 H 是 \(\mathcal{F}(X; Y)\) 的子集。则若 H 关于 \(\mathfrak{S}_1\) -收敛是完全的，它关于 \(\mathfrak{S}_2\) -收敛也是完全的。

因为关于 \(S_{2}\) -收敛的每个柯西滤子也是关于 \(S_{1}\) -收敛的柯西滤子，故可应用推论 1。

推论 3. 设 H 是 \(\mathcal{F}(\mathbf{X};\mathbf{Y})\) 的子集，使得对每个

\[
x \notin \mathrm{B} = \bigcup_ {\mathrm{A} \in \mathfrak {S}} \mathrm{A},
\]

\(\mathbf{H}(x)\) 在 \(\mathbf{Y}\) 中的闭包是 \(\mathbf{Y}\) 的完全子空间。则 \(\overline{\mathbf{H}}\)（即 \(\mathbf{H}\) 在 \(\mathcal{F}_{\mathfrak{S}}(\mathbf{X};\mathbf{Y})\) 中的闭包）是完全子空间。

设 \(\Phi\) 是 \(\overline{\mathbf{H}}\) 上的柯西滤子，如下定义映射 \(v: \mathbf{X} \to \mathbf{Y}\)。若 \(x \in \mathbf{B}\)，则 \(\Phi(x)\) 是 \(\overline{\mathbf{H}(x)}\) 上的柯西滤子（第 2 号注记 6），故由假设它至少有一个极限点；取 \(v(x)\) 为这些极限之一。若 \(x \notin \mathbf{B}\)，取 \(v(x)\) 为 \(\mathbf{Y}\) 的任意一点。以如此定义的 \(v\)，显然 \(\Phi\) 逐点收敛于 \(v\)（在 \(\mathbf{B}\) 中），因而由命题 5，\(v\) 是 \(\Phi\) 在 \(\mathcal{F}_{\mathfrak{S}}(\mathbf{X}; \mathbf{Y})\) 中的一个极限。

特别地，若 Y 是完全的，则命题 5 的推论 3 的假设对每个 \(H \subset \mathcal{F}(X; Y)\) 都满足；因此：

定理 1. 设 X 是一个集，S 是 X 的子集组成的一个集，Y 是完全一致空间。则一致空间 \(\mathcal{F}_{\mathfrak{S}}(X;Y)\) 是完全的。
% semantic-fragments: legacy-input-seam
\relax{}
